\section{Pilot-0.2 Literature Positioning Audit}
Der Audit findet keinen SAME-Treffer für den \emph{vollständigen} integrierten P2-A-Befund, wohl aber starke Nachbarschaften. Foures--Caulfield--Schmid und Marcotte--Caulfield zeigen bereits, dass energy-optimal perturbations für Mixing stark suboptimal sein können. Sévellec et al. demonstrieren objective-dependent optimale Anfangsstörungen für geophysikalische Transportgrößen. Landreman--Plunk--Dorland zeigen transient free-energy growth und particle/heat flux in modally stable Plasma-Systemen, optimieren den Flux aber nicht separat über Anfangsstörungen. Plunk--Helander behandeln optimal free-energy growth, nicht ein paralleles signed particle-flux objective.

Daher ist ausdrücklich \emph{nicht} neu:
\[
\text{energy-optimal}\neq\text{transport-/mixing-optimal}
\]
als allgemeines Prinzip. Der engere, in der geprüften Literatur nicht als SAME gefundene Anwendungsclaim ist:
\[
\boxed{
\begin{gathered}
\text{spectrally stable drift/transport dynamics}\\
+\ \text{finite-time energy amplification}\\
+\ \text{physically signed particle transport}\\
+\ \text{separate optimization of both}\\
+\ \text{large, resolution-robust quantitative gap}.
\end{gathered}}
\]
Der publizistisch stärkste Witness bleibt $\Delta_\Gamma(1)\simeq0.504$.

\section{Cross-Domain Application Architecture}
Nach P2-A dient die zweite/weitere Anwendung nicht mehr dazu, das Framework überhaupt zu retten, sondern seine wissenschaftliche Tragfähigkeit über einen kontrollierten Plasmafall hinaus zu testen. Die aktuelle Architektur ist:
\[
\boxed{
\text{D10-ZF P2-A}
\longrightarrow
\begin{cases}
\text{Neuro} & \text{active, geometry gate}\\
\text{Climate/Ocean} & \text{active, numerical qualification}\\
\text{Power Grids} & \text{protected}\\
\text{Photonics/Waves} & \text{protected}\\
\text{realistic Fusion} & \text{protected domain-depth}
\end{cases}}
\]
